% Part: proof-theory
% Chapter: sequent-calculus
% Section: invertibility

\documentclass[../../../include/open-logic-section]{subfiles}

\begin{document}

\olfileid{pt}{seq}{inv}

\olsection{বিধির উল্টানোযোগ্যতা}

\begin{defn}
কোনো বিধি~$R$,
\[
\AxiomC{$S_1$}
\AxiomC{$\dots$}
\AxiomC{$S_n$}
\RightLabel{$R$}
\TrinaryInfC{$S$}
\DisplayProof
\]
কোনো ব্যবস্থায় \emph{উল্টানোযোগ্য}, যদি
$\Proves S$ হলেই $\Proves S_i$ হয়,
$i = 1$, \dots,~$n$-এর প্রত্যেকটির
জন্য। বিধিটি \emph{!!{height}-সংরক্ষণকারী
উল্টানোযোগ্য} (বা \emph{!!{hp}-উল্টানোযোগ্য}),
যদি $\Proves[n] S$ হলেই
$\Proves[n] S_i$ হয়,
$i = 1$, \dots,~$n$-এর প্রত্যেকটির
জন্য।
\end{defn}
% BN-SRC-629 TARGET-CORRECTION-BEGIN
\par\smallskip\noindent\textnormal{\small\textbf{উৎসসংশোধন (BN-SRC-629):} সাধারণ উল্টানোযোগ্যতার সংজ্ঞায় উৎসের অনির্ধারিত \(h_i\) উচ্চতা বাদ দিয়ে পূর্বধারণার সাধারণ প্রমাণযোগ্যতা রাখা হয়েছে। উচ্চতা-সংরক্ষণকারী সংজ্ঞার উপপক্ষে উৎসের \(\Proves_n S\) বদলে পাশের বিধিবদ্ধ লিখন \(\Proves[n] S\) করা হয়েছে।}
% BN-SRC-629 TARGET-CORRECTION-END

\begin{lem}\ollabel{lem:G3c-invert}
\Log{G3c}-তে $\lnot$, $\land$, $\lor$
ও $\lif$-এর বিধিগুলি
!!{height}-সংরক্ষণকারী উল্টানোযোগ্য;
অর্থাৎ:
\begin{enumerate}
  \item \RightR{\lnot}: যদি $\Log{G3c} \Proves[n] \Gamma \Sequent
  \Delta, \lnot !A$ হয়, তবে $\Log{G3c} \Proves[n] !A, \Gamma \Sequent
  \Delta$.
  \item \LeftR{\lnot}: যদি $\Log{G3c} \Proves[n] \Gamma, \lnot !A
  \Sequent \Delta$ হয়, তবে $\Log{G3c} \Proves[n] \Gamma \Sequent \Delta,
  !A$.
  \item \RightR{\land}: যদি $\Log{G3c} \Proves[n] \Gamma \Sequent
  \Delta, !A \land !B$ হয়, তবে $\Log{G3c} \Proves[n] \Gamma \Sequent
  \Delta, !A$ এবং $\Log{G3c} \Proves[n] \Gamma \Sequent
  \Delta, !B$.
  \item \LeftR{\land}: যদি $\Log{G3c} \Proves[n] \Gamma, !A \land !B
  \Sequent \Delta$ হয়, তবে $\Log{G3c} \Proves[n] !A, !B, \Gamma \Sequent
  \Delta$.
  \item \RightR{\lor}: যদি $\Log{G3c} \Proves[n] \Gamma \Sequent
  \Delta, !A \lor !B$ হয়, তবে $\Log{G3c} \Proves[n] \Gamma \Sequent
  \Delta, !A, !B$.
  \item \LeftR{\lor}: যদি $\Log{G3c} \Proves[n] !A \lor !B, \Gamma \Sequent
  \Delta$ হয়, তবে $\Log{G3c} \Proves[n] !A, \Gamma \Sequent
  \Delta$ এবং $\Log{G3c} \Proves[n] !B, \Gamma \Sequent
  \Delta$.
  \item \RightR{\lif}: যদি $\Log{G3c} \Proves[n] \Gamma \Sequent
  \Delta, !A \lif !B$ হয়, তবে $\Log{G3c} \Proves[n] !A, \Gamma \Sequent
  \Delta, !B$.
  \item \LeftR{\lif}: যদি $\Log{G3c} \Proves[n] !A \lif !B, \Gamma
  \Sequent \Delta$ হয়, তবে $\Log{G3c} \Proves[n] \Gamma \Sequent \Delta,
  !A$ এবং $\Log{G3c} \Proves[n] !B, \Gamma \Sequent \Delta$।
\end{enumerate}
\end{lem}

\begin{proof}
\Log{G3c}-র যেকোনো বিধি~$R$-এর
জন্য দেখাতে হবে: $R$-এর সিদ্ধান্ত~$S$-এর
!!a{proof}~$\pi$ থাকলে বিধি~$R$-এর পূর্বধারণা
$S_i$-গুলিরও !!{proof}s~$\pi_i$
থাকে, যেখানে $\pheight{\pi_i}
\le \pheight{\pi}$। \LeftR{\land}
বিধি বিবেচনা করি। ধরি,
$!A \land !B, \Gamma
\Sequent \Delta$ রূপের অন্তিম
সিকোয়েন্টযুক্ত !!a{proof}~$\pi$
আছে। দেখাতে হবে,
$!A, !B, \Gamma \Sequent \Delta$-র
!!a{proof} $\pi'$ আছে, যাতে
$\pheight{\pi'} \le
\pheight{\pi}$। আমরা
$n = \pheight{\pi}$-এর ওপর
আরোহ করে এটি দেখাই।

$n = 0$ হলে $\pi$ কেবল একটি
স্বতঃসিদ্ধ: $!A \land !B, \Gamma
\Sequent \Delta$। \Log{G3c}-র
স্বতঃসিদ্ধ হয় আণবিক $!C$-সহ
$!C, \Gamma
\Sequent \Delta, !C$ রূপের,
নয়তো বাঁ দিকে অসংগতি-সূত্রসহ
সিকোয়েন্ট। $!A \land !B$
আণবিক নয়, তাই পরিচয়-স্বতঃসিদ্ধে
এটি $!C$ হতে পারে না।
$!A \land !B, \Gamma \Sequent
\Delta$ পরিচয়-স্বতঃসিদ্ধ হলে
কোনো আণবিক $!C$ অবশ্যই
$\Gamma$ ও~$\Delta$ উভয়ের
!!a{element}। আর মিথ্যা-স্বতঃসিদ্ধ
হলে অসংগতি-সূত্রটি প্রসঙ্গে থাকে।
দুই ক্ষেত্রেই $!A, !B, \Gamma
\Sequent \Delta$ স্বতঃসিদ্ধ;
তাই চাওয়া !!{proof}~$\pi'$
পাওয়া যায় (এর
!!{height}ও~$0$)।
% BN-SRC-630 TARGET-CORRECTION-BEGIN
\par\smallskip\noindent\textnormal{\small\textbf{উৎসসংশোধন (BN-SRC-630):} G3c-র প্রদর্শিত বিধি-সারণিতে আণবিক পরিচয়-স্বতঃসিদ্ধের পাশাপাশি বাঁদিকে \(\lfalse\)-সহ স্বতঃসিদ্ধও আছে। শূন্য-উচ্চতার আরোহী ক্ষেত্রে এই দ্বিতীয় স্বতঃসিদ্ধও পরীক্ষা করা হয়েছে; ফল একই থাকে।}
% BN-SRC-630 TARGET-CORRECTION-END

$n>0$ হলে প্রমাণটির শেষে
কোনো অনুমান-ধাপ আছে।
$n$-এর জন্য দাবিটি প্রমাণ করতে হবে,
তবে ভিন্ন প্রসঙ্গ হলেও
উচ্চতা~$<n$-এর যেকোনো
!!{proof}-এর জন্য এটি সত্য
ধরতে পারি। অর্থাৎ যেকোনো
$k<n$ এবং যেকোনো
$\Pi$, $\Lambda$-র ক্ষেত্রে,
$!A \land !B, \Pi \Sequent
\Lambda$-র !!{height}~$k$-এর
!!a{proof} থাকলে
$!A, !B, \Pi \Sequent \Lambda$-র
!!{height}~$\le k$-এর
!!a{proof} আছে।

দুটি ক্ষেত্র: বাঁ দিকের
$!A \land !B$ মুখ্য,
অথবা মুখ্য নয়। মুখ্য হলে
শেষ অনুমান \LeftR{\land},
আর অব্যবহিত
উপ-!!{proof}~$\pi_1$-এর
অন্তিম সিকোয়েন্ট
$!A, !B, \Gamma \Sequent
\Delta$। এখানে
$\pheight{\pi_1} <
\pheight{\pi}$ বলে
দাবিটি প্রতিষ্ঠিত।

$!A \land !B$ মুখ্য না হলে
অন্য কোনো !!{formula} মুখ্য,
আর $!A \land !B$ শেষ
অনুমানের প্রসঙ্গে থাকে।
তখন শেষ বিধি~$R$-এর
অব্যবহিত উপ-!!{proof}s-এর
ওপর আরোহের অনুমান প্রয়োগ করি।
এতে $\pi_1'$ নামে
!!a{proof} অথবা
$\pi_1'$ ও~$\pi_2'$ নামে
দুটি !!{proof}s পাওয়া যায়;
তাদের !!{height}
$\pi$-র সংশ্লিষ্ট অব্যবহিত
উপ-!!{proof}s-এর উচ্চতার
বেশি নয়। $\pi_1'$ ও
$\pi_2'$-এর অন্তিম
সিকোয়েন্টগুলিতে বাঁ দিকের
প্রসঙ্গে $!A \land !B$-এর
জায়গায় $!A, !B$ থাকে।
তারপর বিধি~$R$ প্রয়োগ
করে প্রমাণ~$\pi'$ পাই।
যেমন, ধরি $\pi$-র
শেষে~\LeftR{\lif} আছে:
\[
\AxiomC{}
\RightLabel{$\pi_1$}
\Deduce$!A \land !B, \Gamma' \fCenter \Delta, !C$
\AxiomC{}
\RightLabel{$\pi_2$}
\Deduce$!A \land !B, \Gamma', !D \fCenter \Delta$
\RightLabel{\LeftR{\lif}}
\BinaryInf$!A \land !B, \Gamma', !C \lif !D \fCenter \Delta$
\DisplayProof
\]
এখানে বাঁ দিকের
$!C \lif !D$ মুখ্য,
আর বাঁ দিকের প্রসঙ্গ
$!A \land !B, \Gamma'$
(অর্থাৎ $\Gamma = \Gamma',
!C \lif !D$)। আরোহের
অনুমানে যথাক্রমে
$!A, !B, \Gamma' \Sequent
\Delta, !C$ এবং
$!A, !B, \Gamma', !D
\Sequent \Delta$-র
!!{proof}s $\pi_1'$ ও
$\pi_2'$ পাওয়া যায়।
এই দুটি সিকোয়েন্টকে
\LeftR{\lif}-এর পূর্বধারণা
করে~$\pi'$ পাওয়া যায়:
\[
\AxiomC{}
\RightLabel{$\pi_1'$}
\Deduce$!A, !B, \Gamma' \fCenter \Delta, !C$
\AxiomC{}
\RightLabel{$\pi_2'$}
\Deduce$!A, !B, \Gamma', !D \fCenter \Delta$
\RightLabel{\LeftR{\lif}}
\BinaryInf$!A, !B, \Gamma', !C \lif !D \fCenter \Delta$
\DisplayProof
\]
এখন,
\begin{align*}
\pheight{\pi} & = \max(\pheight{\pi_1}, \pheight{\pi_2}) + 1\\
\pheight{\pi'} & = \max(\pheight{\pi_1'}, \pheight{\pi_2'}) + 1
\intertext{সংজ্ঞা অনুসারে, আর}
\pheight{\pi_1'} & \le \pheight{\pi_1}\\
\pheight{\pi_2'} & \le \pheight{\pi_2}
\intertext{আরোহের অনুমান অনুসারে। অতএব,}
\pheight{\pi'} & \le \pheight{\pi}.
\end{align*}
\end{proof}
% BN-SRC-631 TARGET-CORRECTION-BEGIN
\par\smallskip\noindent\textnormal{\small\textbf{উৎসসংশোধন (BN-SRC-631):} \(\LeftR{\lif}\) উদাহরণে আরোহের অনুমান থেকে পাওয়া \(\pi'_1,\pi'_2\)-র অন্তিম সিকোয়েন্টে \(A\land B\)-এর বদলে \(A,B\) থাকা চাই। উৎসের গদ্যে পুরোনো প্রসঙ্গই পুনরাবৃত্ত, অথচ পরের প্রদর্শিত বৃক্ষে সঠিক \(A,B\) আছে; গদ্যটি বৃক্ষের সঙ্গে মেলানো হয়েছে।}
% BN-SRC-631 TARGET-CORRECTION-END
% BN-NORM-181 TARGET-NORMALIZATION-BEGIN
\par\smallskip\noindent\textnormal{\small\textbf{ভাষাগত স্বাভাবিকীকরণ (BN-NORM-181):} উচ্চতা-সমীকরণের মধ্যবর্তী দুটি ইংরেজি বাক্য বাংলা করা হয়েছে; সব সমীকরণ ও তুলনাচিহ্ন অপরিবর্তিত।}
% BN-NORM-181 TARGET-NORMALIZATION-END

\begin{prob}
\RightR{\lif}-এর ক্ষেত্রে
\olref{lem:G3c-invert}-এর
প্রমাণ দাও।
\end{prob}

\begin{lem}\ollabel{lem:invert-quant}
\RightR{\lforall} ও
\LeftR{\lexists} বিধি
নিচের অর্থে
!!{height}-সংরক্ষণকারী উল্টানোযোগ্য:
\begin{enumerate}
  \item যদি $\Log{G3c} \Proves[n] \Gamma \Sequent \Delta,
  \lforall[x][!A(x)]$ হয়, তবে $\Log{G3c} \Proves[n] \Gamma \Sequent
  \Delta, !A(c)$; এবং
  \item যদি $\Log{G3c} \Proves[n] \lexists[x][!A(x)], \Gamma \Sequent
  \Delta$ হয়, তবে $\Log{G3c} \Proves[n] !A(c), \Gamma \Sequent \Delta$,
\end{enumerate}
যেকোনো $c$-র জন্য,
যা যথাক্রমে $\Gamma$,
$\Delta$ এবং
$\lforall[x][!A(x)]$
অথবা $\lexists[x][!A(x)]$-এ
উপস্থিত নয়।
\end{lem}

\begin{proof}
  (১) প্রমাণ করি; (২)
  অনুশীলন হিসেবে থাক।
  যুক্তিটি \olref{lem:G3c-invert}-এর
  প্রমাণের মতো। আরোহী ধাপে
  আবার দুটি ক্ষেত্র আছে।
  $\lforall[x][!A(x)]$
  মুখ্য হলে কাজটি সংক্ষিপ্ত:
  শেষ $\RightR{\lforall}$
  অনুমানের পূর্বধারণা চাওয়া
  $\Gamma \Sequent \Delta, !A(a)$
  রূপের, আর সেখানে শেষ হওয়া
  অব্যবহিত উপ-!!{proof}~$\pi_1$-এর
  $\pheight{\pi_1} <
  \pheight{\pi}$। অধিকন্তু
  আইগেনচলের শর্তে $a$,
  $\Gamma$, $\Delta$ কিংবা
  $\lforall[x][!A(x)]$-এ
  উপস্থিত নয়। $\pi_1$
  নিয়মিত ধরা যায়; তাই
  \olref[qua]{lem:subst-regular}
  অনুসারে $\Subst{\pi_1}{c}{a}$
  একই !!{height}-এর
  $\Gamma \Sequent \Delta, !A(c)$-র
  !!a{proof}।

  দ্বিতীয় ক্ষেত্রে
  $\lforall[x][!A(x)]$
  শেষ অনুমানে মুখ্য নয়;
  অন্য কোনো সূত্র মুখ্য।
  আবার একটি উদাহরণ দেখি।
  ধরি, শেষ অনুমান
  \RightR{\land}।
  $\Delta$-র কোনো
  !!{formula}~$!B \land !C$
  রূপের এবং মুখ্য।
  !!{proof}~$\pi$-র শেষ অংশ:
  \[
\AxiomC{}
\RightLabel{$\pi_1$}
\Deduce$\Gamma \fCenter \Delta', \lforall[x][!A(x)], !B$
\AxiomC{}
\RightLabel{$\pi_2$}
\Deduce$\Gamma \fCenter \Delta', \lforall[x][!A(x)], !C$
\RightLabel{\RightR{\land}}
\BinaryInf$\Gamma \fCenter \Delta', \lforall[x][!A(x)], !B \land !C$
\DisplayProof
\]
এখানে $\Delta = \Delta',
!B \land !C$। উপপাদ্যের
সম্পূর্ণ সতেজতা-শর্ত
মেনে $c$ এমন
!!a{constant}, যা
$\Delta$-তে নেই।
বিশেষত এটি
$\Delta', !B$ বা
$\Delta', !C$-তেও
নেই। তাই $\pi_1$
ও~$\pi_2$-তে আরোহের
অনুমান প্রযোজ্য;
!!{proof}s~$\pi_1'$ ও
$\pi_2'$ পাওয়া যায়,
যেখানে $\pheight{\pi_1'}
\le \pheight{\pi_1}$
এবং $\pheight{\pi_2'}
\le \pheight{\pi_2}$।
$\Gamma \Sequent
\Delta, !A(c)$-র
চাওয়া !!{proof}~$\pi'$:
\[
\AxiomC{}
\RightLabel{$\pi_1'$}
\Deduce$\Gamma \fCenter \Delta', !A(c), !B$
\AxiomC{}
\RightLabel{$\pi_2'$}
\Deduce$\Gamma \fCenter \Delta', !A(c), !C$
\RightLabel{\RightR{\land}}
\BinaryInf$\Gamma \fCenter \Delta', !A(c), !B \land !C$
\DisplayProof
\]
এখানে $\pheight{\pi'} = \max(\pheight{\pi_1'}, \pheight{\pi_2'})+1 \le \pheight{\pi}$।
\end{proof}
% BN-SRC-632 TARGET-CORRECTION-BEGIN
\par\smallskip\noindent\textnormal{\small\textbf{উৎসসংশোধন (BN-SRC-632):} অমুখ্য \(\RightR{\land}\) ক্ষেত্রে উৎস কেবল \(\Delta\)-তে \(c\) অনুপস্থিত বলে; পরিমাণসূচক-বিধির দাবির জন্য \(c\) অবশ্যই \(\Gamma\), \(\Delta\) ও \(\lforall[x][!A(x)]\) থেকে নতুন হতে হবে। উদাহরণেও পূর্ণ সতেজতা-শর্ত স্পষ্ট করা হয়েছে।}
% BN-SRC-632 TARGET-CORRECTION-END

\Olref{lem:G3c-invert}
কর্তন-বর্জন উপপাদ্যের
প্রমাণে গুরুত্বপূর্ণ।
এটি দিয়ে নিচের
ফলটিও দেখানো যায়:

\begin{prop}\ollabel{prop:G3c-cont-adm}
\Log{G1c}-র
\LeftR{\Contraction}
ও \RightR{\Contraction}
সংকোচন-বিধি \Log{G3c}-তে
!!{height}-সংরক্ষণকারী
অনুমোদনযোগ্য।
\end{prop}

\begin{proof}
\RightR{\Contraction}
ও \LeftR{\Contraction}-এর
যুক্তি একসঙ্গে দিই।
ধরি, \Log{G3c}-তে
$\Gamma \Sequent \Delta, !A, !A$
অথবা $!A, !A, \Gamma
\Sequent \Delta$-র
!!a{proof} হলো $\pi$।
দেখাতে হবে,
যথাক্রমে $\Gamma
\Sequent \Delta, !A$
অথবা $!A, \Gamma
\Sequent \Delta$-রও
\Log{G3c}-!!{proof}~$\pi'$
আছে, যেখানে
$\pheight{\pi'} \le
\pheight{\pi}$।
$n = \pheight{\pi}$-এর
ওপর আরোহ করি।

$n=0$ হলে $\pi$
কেবল একটি স্বতঃসিদ্ধ।
$\Gamma \Sequent \Delta, !A,
!A$ যদি \Log{G3c}-র
স্বতঃসিদ্ধ হয়, তবে
$\Gamma \Sequent \Delta, !A$-ও
স্বতঃসিদ্ধ। কারণ হয় কোনো
আণবিক~$!C$, $\Gamma$
ও~$\Delta$ উভয়ের
!!a{element}; নয়তো
$!A$ আণবিক এবং
$\Gamma$-র !!a{element}।
এছাড়া বাঁ প্রসঙ্গে
অসংগতি-সূত্র থাকলেও
দুটি সিকোয়েন্টই
স্বতঃসিদ্ধ। অন্তিম
সিকোয়েন্ট $!A, !A,
\Gamma \Sequent \Delta$
হলেও একই যুক্তি চলে।
% BN-SRC-633 TARGET-CORRECTION-BEGIN
\par\smallskip\noindent\textnormal{\small\textbf{উৎসসংশোধন (BN-SRC-633):} সংকোচনের \(n=0\) ক্ষেত্রে G3c-র আণবিক পরিচয়-স্বতঃসিদ্ধ ছাড়াও বাঁদিকে \(\lfalse\)-সহ স্বতঃসিদ্ধ আছে। সংকুচিত সিকোয়েন্টে সেই কারণও অক্ষত থাকে; উৎসে অনুলিখিত ক্ষেত্রটি যোগ করা হয়েছে।}
% BN-SRC-633 TARGET-CORRECTION-END

$n>0$ হলে $\pi$-র
শেষে কোনো বিধি~$R$
আছে। এই শেষ অনুমানে
$!A$ মুখ্য না হলে
\olref{lem:G3c-invert}-এর
প্রমাণের মতো $\pi$-র
অব্যবহিত উপ-!!{proof}s-এ
আরোহের অনুমান প্রয়োগ
করে, তার ফলের
!!{proof}(s)-এ একই
বিধি~$R$ প্রয়োগ করি।
এতে !!a{proof}~$\pi'$
পাই। আরোহের অনুমান
হলো: $\Pi \Sequent
\Lambda, !D, !D$
অথবা $!D, !D, \Pi
\Sequent \Lambda$-র
!!{height}~$k<n$-এর
!!a{proof} থাকলে
যথাক্রমে $\Pi
\Sequent \Lambda, !D$
অথবা $!D, \Pi
\Sequent \Lambda$-র
!!{height}~$\le k$-এর
!!a{proof} আছে।
(প্রসঙ্গ বা সংকুচিত
!!{formula}, $\Gamma$,
$\Delta$ বা~$!A$
থেকে আলাদা হলেও
আরোহের অনুমান প্রযোজ্য।)

যেমন, ধরি $R$ হলো
$\RightR{\land}$ এবং
$\pi$-র শেষ অংশ:
\[
\AxiomC{}
\RightLabel{$\pi_1$}
\Deduce$\Gamma \fCenter \Delta', !B, !A, !A$
\AxiomC{}
\RightLabel{$\pi_2$}
\Deduce$\Gamma \fCenter \Delta', !C, !A, !A$
\RightLabel{\RightR{\land}}
\BinaryInf$\Gamma \fCenter \Delta', !B \land !C, !A, !A$
\DisplayProof
\]
এখানে $\Delta = \Delta',
!B \land !C$। আরোহের
অনুমানে যথাক্রমে
$\Gamma \fCenter
\Delta', !B, !A$
ও $\Gamma \fCenter
\Delta', !C, !A$-র
!!{proof}s $\pi_1'$
ও $\pi_2'$ পাওয়া যায়;
$\pheight{\pi_1'} \le
\pheight{\pi_1}$
এবং $\pheight{\pi_2'}
\le \pheight{\pi_2}$।
চাওয়া !!{proof}~$\pi'$:
\[
\AxiomC{}
\RightLabel{$\pi_1'$}
\Deduce$\Gamma \fCenter \Delta', !B, !A$
\AxiomC{}
\RightLabel{$\pi_2'$}
\Deduce$\Gamma \fCenter \Delta', !C, !A$
\RightLabel{\RightR{\land}}
\BinaryInf$\Gamma \fCenter \Delta', !B \land !C, !A$
\DisplayProof
\]

শেষ অনুমানে $!A$
মুখ্য হলে $\pi$-র
অব্যবহিত উপ-!!{proof}s-এ
প্রথমে উল্টানোযোগ্যতার
সহায়ক উপপাদ্য
(\olref{lem:G3c-invert}),
তারপর আরোহের অনুমান,
শেষে আবার বিধি~$R$
প্রয়োগ করি। যেমন,
ধরি $\pi$-র সিদ্ধান্ত
$\Gamma \Sequent
\Delta, !A, !A$,
$!A \ident (!B \land !C)$,
এবং শেষ অনুমানে সূত্রটি
মুখ্য। তাহলে $\pi$-র
রূপ:
\[
\AxiomC{}
\RightLabel{$\pi_1$}
\Deduce$\Gamma \fCenter \Delta, !B \land !C, !B$
\AxiomC{}
\RightLabel{$\pi_2$}
\Deduce$\Gamma \fCenter \Delta, !B \land !C, !C$
\RightLabel{\RightR{\land}}
\BinaryInf$\Gamma \fCenter \Delta, !B \land !C, !B \land !C$
\DisplayProof
\]
\olref{lem:G3c-invert}
$\pi_1$-এ প্রয়োগ করলে
$\Gamma \Sequent
\Delta, !B, !B$-র
!!a{proof}~$\pi_1'$
পাই। (এতে $\Gamma
\Sequent \Delta,
!C, !B$-রও
!!a{proof} মেলে,
কিন্তু সেটি লাগবে না।)
একইভাবে $\pi_2$-এ
\olref{lem:G3c-invert}
প্রয়োগ করে $\Gamma
\Sequent \Delta,
!C, !C$-র
!!a{proof}~$\pi_2'$
পাই। \RightR{\land}
উল্টানো !!{height}-সংরক্ষণকারী
বলে $\pheight{\pi_1'}
\le \pheight{\pi_1}
< \pheight{\pi}$ এবং
$\pheight{\pi_2'}
\le \pheight{\pi_2}
< \pheight{\pi}$।
তাই $\pi_1'$ ও
$\pi_2'$-এ আরোহের
অনুমান প্রযোজ্য:
যথাক্রমে $\Gamma
\Sequent \Delta, !B$
ও $\Gamma \Sequent
\Delta, !C$-র
!!{proof}s~$\pi_1''$
ও $\pi_2''$ আছে,
যেখানে $\pheight{\pi_1''}
\le \pheight{\pi_1'}
\le \pheight{\pi_1}$
এবং $\pheight{\pi_2''}
\le \pheight{\pi_2'}
\le \pheight{\pi_2}$।
তখন !!{proof}~$\pi'$:
\[
\AxiomC{}
\RightLabel{$\pi_1''$}
\Deduce$\Gamma \fCenter \Delta, !B$
\AxiomC{}
\RightLabel{$\pi_2''$}
\Deduce$\Gamma \fCenter \Delta, !C$
\RightLabel{\RightR{\land}}
\BinaryInf$\Gamma \fCenter \Delta, !B \land !C$
\DisplayProof
\]
এতে $\pheight{\pi'}
\le \pheight{\pi}$।

!!{formula} মুখ্য হলে
পরিমাণসূচক-বিধিগুলির
ক্ষেত্রে বিশেষ যত্ন দরকার।

ধরি, ডান দিকে
$!A \ident \lforall[x][!B(x)]$
মুখ্য। তাহলে
$\pi$-র শেষ অংশ:
\[
\AxiomC{}
\RightLabel{$\pi_1$}
\Deduce$\Gamma \fCenter \Delta, \lforall[x][!B(x)], !B(a)$
\RightLabel{\RightR{\lforall}}
\UnaryInf$\Gamma \fCenter \Delta, \lforall[x][!B(x)], \lforall[x][!B(x)]$
\DisplayProof
\]
\olref{lem:invert-quant}
অনুসারে $\Gamma \fCenter
\Delta, !B(c), !B(a)$-র
!!a{proof}~$\pi_1'$
আছে, যার
!!{height} $\le
\pheight{\pi_1}$।
এটি সেই যেকোনো~$c$-র
জন্য, যা $\Gamma
\fCenter \Delta,
\lforall[x][!B(x)], !B(a)$-তে
উপস্থিত নয়। $\pi_1'$
নিয়মিত ধরা যায়;
তাই \olref[qua]{lem:subst-regular}
অনুসারে !!{proof}
$\Subst{\pi_1'}{a}{c}$
হলো $\Gamma \fCenter
\Delta, !B(a), !B(a)$-র
!!a{proof}, যার
!!{height}ও $\le
\pheight{\pi_1}$।
এখন আরোহের অনুমান
থেকে $\Gamma
\fCenter \Delta, !B(a)$-র
!!a{proof}~$\pi_1''$
পাই। \RightR{\lforall}
প্রয়োগ করে~$\pi'$:
\[
\AxiomC{}
\RightLabel{$\pi_1''$}
\Deduce$\Gamma \fCenter \Delta, !B(a)$
\RightLabel{\RightR{\lforall}}
\UnaryInf$\Gamma \fCenter \Delta, \lforall[x][!B(x)]$
\DisplayProof
\]
এখানে $\pheight{\pi'}
\le \pheight{\pi}$।
% BN-SRC-634 TARGET-CORRECTION-BEGIN
\par\smallskip\noindent\textnormal{\small\textbf{উৎসসংশোধন (BN-SRC-634):} \(\lforall\)-ডান মুখ্য সংকোচনের প্রথম বৃক্ষে উৎসে ভুল করে \(\RightR{\lexists}\) লেবেল আছে। পূর্বধারণা, সিদ্ধান্ত এবং পরের প্রয়োগ অনুযায়ী লেবেল \(\RightR{\lforall}\) করা হয়েছে।}
% BN-SRC-634 TARGET-CORRECTION-END

এবার ধরি, $!A \ident
\lexists[x][!B(x)]$
ডান দিকে মুখ্য।
তাহলে $\pi$-র রূপ:
\[
\AxiomC{}
\RightLabel{$\pi_1$}
\Deduce$\Gamma \fCenter \Delta, \lexists[x][!B(x)], \lexists[x][!B(x)], !B(t)$
\RightLabel{\RightR{\lexists}}
\UnaryInf$\Gamma \fCenter \Delta, \lexists[x][!B(x)], \lexists[x][!B(x)]$
\DisplayProof
\]
$\pi_1$-এ আরোহের অনুমান
প্রযোজ্য। ফলে $\Gamma
\Sequent \Delta,
\lexists[x][!B(x)], !B(t)$-র
!!a{proof}~$\pi_1'$ আছে।
(এখানে উল্টানোযোগ্যতার
সহায়ক উপপাদ্য লাগে না।)
তখন !!{proof}~$\pi'$:
\[
\AxiomC{}
\RightLabel{$\pi_1'$}
\Deduce$\Gamma \fCenter \Delta,  \lexists[x][!B(x)], !B(t)$
\RightLabel{\RightR{\lexists}}
\UnaryInf$\Gamma \fCenter \Delta, \lexists[x][!B(x)]$
\DisplayProof
\]
এতে $\pheight{\pi'} \le
\pheight{\pi}$।
\end{proof}
% BN-SRC-635 TARGET-CORRECTION-BEGIN
\par\smallskip\noindent\textnormal{\small\textbf{উৎসসংশোধন (BN-SRC-635):} \(\lexists\)-ডান মুখ্য সংকোচনে উৎসের শেষ গদ্য ও প্রমাণবৃক্ষে সাক্ষী-পদ \(t\) বাদ পড়ে \(B\) লেখা। আগের পূর্বধারণার সঙ্গে মিলিয়ে \(B(t)\) রাখা হয়েছে। একবচন \(\pi'_1\)-এর জন্য উৎসের ভুল বহুবচন proof-token-ও একবচন করা হয়েছে।}
% BN-SRC-635 TARGET-CORRECTION-END

\begin{prob}
\LeftR{\Contraction}-এর
ক্ষেত্রে \olref{prop:G3c-cont-adm}-এর
প্রমাণের রূপরেখা দাও।
$!A \ident (!B \land !C)$
এবং $!A \ident (!B
\lif !C)$ ক্ষেত্রগুলি
বর্ণনা করো।
\end{prob}

\begin{prob}
অবশিষ্ট পরিমাণসূচক-ক্ষেত্রগুলিতে
\LeftR{\Contraction}-এর
জন্য \olref{prop:G3c-cont-adm}-এর
প্রমাণের রূপরেখা দাও;
অর্থাৎ $!A \ident
\lforall[x][!B(x)]$ এবং
$!A \ident \lexists[x][!B(x)]$
ক্ষেত্রে।
\end{prob}


\end{document}
